\begin{abstract}
Minimum-mass design of steel frames is usually carried out against deterministic code checks,
which leaves the reliability of the optimised frame unquantified. This paper presents an open,
reproducible framework for reliability-constrained, multi-objective optimisation of planar steel
moment frames built on the PyNite finite element library. Discrete W-shape sections are checked
against AISC 360-16 with the Direct Analysis Method. The ultimate limit state is defined at
component level: it is the first attainment of a member strength limit state in a second-order
elastic analysis ($B_1$--$B_2$ amplification of superposed linear responses), not structural
collapse. Fifty-year drift and deflection serviceability are treated alongside it under load,
material, fabrication and imperfection uncertainty, with model-form uncertainty examined
separately. A sample average approximation with common random numbers turns the reliability
constraints into deterministic ones for a genetic algorithm and NSGA-II. Every optimum is
re-verified with $10^6$ independent samples and reported with exact binomial confidence
intervals, and a code-compliant variant requires the deterministic AISC and drift checks in
addition to the reliability targets. The surrogate is validated against OpenSeesPy and against
full second-order re-analysis for eight designs, including targeted samples near the limit
state; reliability indices differ by at most 0.07. For seven frames that vary the number of
storeys, the bay width and the wind pressure, the code-compliant reliability-based optima met
targets of 3.0 (ultimate) and 1.5 (serviceability) and were 1.1--10.1\% heavier than the
deterministic optima, mainly because the deterministic optima had low drift reliability. Where a deterministic design with a tightened
drift limit (up to $h/600$) met the same targets, the reliability-based optimum was within
$-5.0$ to $+0.6$\% of the lightest such design. Uncertainty-aware optimisation did not reduce the steel mass relative
to code-based optimisation. Its benefit is a least-mass design at a stated and uniform
reliability.
\end{abstract}
